\documentclass[11pt]{amsart}
\usepackage{geometry}                % See geometry.pdf to learn the layout options. There are lots.
\geometry{letterpaper}                   % ... or a4paper or a5paper or ... 
%\geometry{landscape}                % Activate for for rotated page geometry
%\usepackage[parfill]{parskip}    % Activate to begin paragraphs with an empty line rather than an indent
\usepackage{graphicx}
\usepackage{amssymb}
\usepackage{epstopdf}
\DeclareGraphicsRule{.tif}{png}{.png}{`convert #1 `dirname #1`/`basename #1 .tif`.png}

\title{PS 4.12}
%\author{The Author}
%\date{}                                           % Activate to display a given date or no date

\begin{document}
\maketitle
%\section{}
%\subsection{}
\noindent En una papelería se maneja información sobre los 8 modelos diferentes de cuadernos que venden. Se conocen los precios de cada modelo:\\

\begin{center}
$P _{l}$, $P_{2}$, $P_{3}$,  • • ,$P_{8}$
\end{center}

\noindent Donde:

\begin{center}
$P_{i}$ es una variable de tipo real que representa el precio del modelo i.\\

Por otra parte se tiene información sobre las ventas realizadas durante los últimos 30 días. Los datos se presentan de esta forma:\\

$DIA _{1}$, $MOD_{1}$, $CANT_{1}$\\
$DIA_{2}$, $MOD_{2}$, $CANT_{2}$\\
...\\
—1, —1, —1
\end{center}

\noindent Donde:
\begin{quote}
$DIA_{i}$ es una variable de tipo entero que representa el día en que se efectuó la venta i ($1 \leq DLA \leq 30$).\\
$MOD_{i}$ es una variable de tipo entero que representa el modelo que se vendió en la venta i ($1 \leq MODELO \leq 8$).\\
$CANT_{i}$ es una variable de tipo entero que representa el número de cuadernos que se vendió en la venta i.\\
\end{quote}


\noindent Construya un diagrama de flujo que calcule lo siguiente:\\


a) El total recaudado por modelo en los 30 días.\\
b) El total recaudado por día.\\
c) ¿Cuál fue el modelo que más dinero produjo en los 30 días?
\end{document}  